% Options for packages loaded elsewhere
\PassOptionsToPackage{unicode}{hyperref}
\PassOptionsToPackage{hyphens}{url}
\PassOptionsToPackage{dvipsnames,svgnames,x11names}{xcolor}
%
\documentclass[
  12pt,
]{article}

\usepackage{amsmath,amssymb}
\usepackage{setspace}
\usepackage{iftex}
\ifPDFTeX
  \usepackage[T1]{fontenc}
  \usepackage[utf8]{inputenc}
  \usepackage{textcomp} % provide euro and other symbols
\else % if luatex or xetex
  \usepackage{unicode-math}
  \defaultfontfeatures{Scale=MatchLowercase}
  \defaultfontfeatures[\rmfamily]{Ligatures=TeX,Scale=1}
\fi
\usepackage{lmodern}
\ifPDFTeX\else  
    % xetex/luatex font selection
\fi
% Use upquote if available, for straight quotes in verbatim environments
\IfFileExists{upquote.sty}{\usepackage{upquote}}{}
\IfFileExists{microtype.sty}{% use microtype if available
  \usepackage[]{microtype}
  \UseMicrotypeSet[protrusion]{basicmath} % disable protrusion for tt fonts
}{}
\makeatletter
\@ifundefined{KOMAClassName}{% if non-KOMA class
  \IfFileExists{parskip.sty}{%
    \usepackage{parskip}
  }{% else
    \setlength{\parindent}{0pt}
    \setlength{\parskip}{6pt plus 2pt minus 1pt}}
}{% if KOMA class
  \KOMAoptions{parskip=half}}
\makeatother
\usepackage{xcolor}
\usepackage[margin=1in]{geometry}
\setlength{\emergencystretch}{3em} % prevent overfull lines
\setcounter{secnumdepth}{5}
% Make \paragraph and \subparagraph free-standing
\ifx\paragraph\undefined\else
  \let\oldparagraph\paragraph
  \renewcommand{\paragraph}[1]{\oldparagraph{#1}\mbox{}}
\fi
\ifx\subparagraph\undefined\else
  \let\oldsubparagraph\subparagraph
  \renewcommand{\subparagraph}[1]{\oldsubparagraph{#1}\mbox{}}
\fi


\providecommand{\tightlist}{%
  \setlength{\itemsep}{0pt}\setlength{\parskip}{0pt}}\usepackage{longtable,booktabs,array}
\usepackage{calc} % for calculating minipage widths
% Correct order of tables after \paragraph or \subparagraph
\usepackage{etoolbox}
\makeatletter
\patchcmd\longtable{\par}{\if@noskipsec\mbox{}\fi\par}{}{}
\makeatother
% Allow footnotes in longtable head/foot
\IfFileExists{footnotehyper.sty}{\usepackage{footnotehyper}}{\usepackage{footnote}}
\makesavenoteenv{longtable}
\usepackage{graphicx}
\makeatletter
\def\maxwidth{\ifdim\Gin@nat@width>\linewidth\linewidth\else\Gin@nat@width\fi}
\def\maxheight{\ifdim\Gin@nat@height>\textheight\textheight\else\Gin@nat@height\fi}
\makeatother
% Scale images if necessary, so that they will not overflow the page
% margins by default, and it is still possible to overwrite the defaults
% using explicit options in \includegraphics[width, height, ...]{}
\setkeys{Gin}{width=\maxwidth,height=\maxheight,keepaspectratio}
% Set default figure placement to htbp
\makeatletter
\def\fps@figure{htbp}
\makeatother
% definitions for citeproc citations
\NewDocumentCommand\citeproctext{}{}
\NewDocumentCommand\citeproc{mm}{%
  \begingroup\def\citeproctext{#2}\cite{#1}\endgroup}
\makeatletter
 % allow citations to break across lines
 \let\@cite@ofmt\@firstofone
 % avoid brackets around text for \cite:
 \def\@biblabel#1{}
 \def\@cite#1#2{{#1\if@tempswa , #2\fi}}
\makeatother
\newlength{\cslhangindent}
\setlength{\cslhangindent}{1.5em}
\newlength{\csllabelwidth}
\setlength{\csllabelwidth}{3em}
\newenvironment{CSLReferences}[2] % #1 hanging-indent, #2 entry-spacing
 {\begin{list}{}{%
  \setlength{\itemindent}{0pt}
  \setlength{\leftmargin}{0pt}
  \setlength{\parsep}{0pt}
  % turn on hanging indent if param 1 is 1
  \ifodd #1
   \setlength{\leftmargin}{\cslhangindent}
   \setlength{\itemindent}{-1\cslhangindent}
  \fi
  % set entry spacing
  \setlength{\itemsep}{#2\baselineskip}}}
 {\end{list}}
\usepackage{calc}
\newcommand{\CSLBlock}[1]{\hfill\break\parbox[t]{\linewidth}{\strut\ignorespaces#1\strut}}
\newcommand{\CSLLeftMargin}[1]{\parbox[t]{\csllabelwidth}{\strut#1\strut}}
\newcommand{\CSLRightInline}[1]{\parbox[t]{\linewidth - \csllabelwidth}{\strut#1\strut}}
\newcommand{\CSLIndent}[1]{\hspace{\cslhangindent}#1}

\makeatletter
\@ifpackageloaded{caption}{}{\usepackage{caption}}
\AtBeginDocument{%
\ifdefined\contentsname
  \renewcommand*\contentsname{Table of contents}
\else
  \newcommand\contentsname{Table of contents}
\fi
\ifdefined\listfigurename
  \renewcommand*\listfigurename{List of Figures}
\else
  \newcommand\listfigurename{List of Figures}
\fi
\ifdefined\listtablename
  \renewcommand*\listtablename{List of Tables}
\else
  \newcommand\listtablename{List of Tables}
\fi
\ifdefined\figurename
  \renewcommand*\figurename{Figure}
\else
  \newcommand\figurename{Figure}
\fi
\ifdefined\tablename
  \renewcommand*\tablename{Table}
\else
  \newcommand\tablename{Table}
\fi
}
\@ifpackageloaded{float}{}{\usepackage{float}}
\floatstyle{ruled}
\@ifundefined{c@chapter}{\newfloat{codelisting}{h}{lop}}{\newfloat{codelisting}{h}{lop}[chapter]}
\floatname{codelisting}{Listing}
\newcommand*\listoflistings{\listof{codelisting}{List of Listings}}
\makeatother
\makeatletter
\makeatother
\makeatletter
\@ifpackageloaded{caption}{}{\usepackage{caption}}
\@ifpackageloaded{subcaption}{}{\usepackage{subcaption}}
\makeatother
\ifLuaTeX
  \usepackage{selnolig}  % disable illegal ligatures
\fi
\usepackage{bookmark}

\IfFileExists{xurl.sty}{\usepackage{xurl}}{} % add URL line breaks if available
\urlstyle{same} % disable monospaced font for URLs
\hypersetup{
  pdftitle={AI Governance in Warfare: Critical Junctures and International Institutional Responses to the Russia-Ukraine War},
  pdfauthor={Wilson Wong},
  colorlinks=true,
  linkcolor={blue},
  filecolor={Maroon},
  citecolor={Blue},
  urlcolor={Blue},
  pdfcreator={LaTeX via pandoc}}

\title{AI Governance in Warfare: Critical Junctures and International
Institutional Responses to the Russia-Ukraine War}
\author{Wilson Wong}
\date{Invalid Date}

\begin{document}
\maketitle

\renewcommand*\contentsname{Table of Contents}
{
\hypersetup{linkcolor=}
\setcounter{tocdepth}{3}
\tableofcontents
}
\setstretch{2}
\newpage
\begin{center}
\vspace*{2in}
{\Large\textbf{AI Governance in Warfare}}\\
\vspace{0.5cm}
{\large Critical Junctures and International Institutional Responses to the Russia-Ukraine War}\\
\vspace{1.5cm}
{\large Wilson Wong}\\
\vspace{1cm}
Governance of AI in Conflict\\
Dr. Julian Campisi, Dr. Titilayo Soremi\\
POLD02 - Independent Senior Research\\
\vspace{1cm}
April 6th, 2026
\end{center}
\newpage

\section{Abstract}\label{abstract}

Artificial Intelligence (AI) has rapidly integrated into our world of
conflict by reshaping how warfare works, targeting systems, and the
deployment of military technology. The 2022 Russia-Ukraine war
highlighted significant challenges regarding both weapons systems and
the legal and governance frameworks governing AI use, raising concerns
about accountability, the rules of war, and the adequacy of current
international regulatory frameworks. Despite the growing usage of AI in
military contexts, there remains limited empirical analysis of whether
international institutions have meaningfully adapted their governance
responses considering this conflict. This study will look towards the
Ukraine-Russia war as a critical juncture in the global governance of
military AI. Drawing on theories such as global governance and
historical institutionalism, this research analyzes how international
institutions respond to major security shocks and technological
disruptions, especially regarding the distinction between hard-law and
soft-law regulatory approaches. The study will use a qualitative
research design based on document and thematic analysis of official
government documents, reports, and statements issued by international
institutions between 2015-2025. It will test the hypothesis that
intergovernmental organizations produce a greater number of AI
governance outputs in the two years after the invasion started. By
examining the content of these outputs, this research assesses whether
institutional responses represent incremental adaptation or more
substantial governance change. The findings aim to contribute to broader
debates on AI governance, global security, and institutional change in
the context of modern warfare.

\section{Introduction}\label{introduction}

Artificial Intelligence (AI) has become one of the most rapidly
advancing technologies of the 21st century. It has reshaped industries,
governance systems and most significantly the norms of armed conflict.
Autonomous targeting systems, battlefield intelligence platforms,
drones, cyber operations have one element that has improved
significantly because AI has increasingly become standard features of
modern warfare.

The ongoing conflict Russia-Ukraine war with the Russian invasion of
February 24, 2022 has placed these developments at the forefront of
international political debate. The conflict has been described as the
first large-scale conflict in which AI-enabled technologies play a very
key role in military operations, raising urgent questions about
accountability, protection of civilians, and the adequacy of current
international regulatory frameworks (Rakhmetov and Murzagulova
(2025);Kostenko (2022)).

However, as the war progresses the global governance of military AI
remains fragmented and largely non-binding. This governance gap poses
serious risks to international peace and security (Tallberg et al.
(2023);Sweijs and Romansky (2024)). This study asks a direct empirical
question: Has the Russia-Ukraine war constituted a critical juncture in
the global governance of Military AI? Drawing on Historical
Institutionalism (HI) and Global Governance Theory, this research
examines the connection of how international institutions have responded
to the invasion, tests whether a measurable surge in governance outputs
occurred, and analyzes whether the outputs represent substantive
transformative change or if it is only an incremental adaptation within
existing legal frameworks.

This study will make three specific contributions to the literature.
First, it provides empirical evidence through a dataset of 29 official
IGO documents spanning 2015-2025 on whether and how institutional
responses to the invasion changed. Second, it applies Historical
Institutionalism critical juncture framework to the domain of military
AI governance, which remains very underexplored in the existing
scholarship. Lastly, it uses the 1997 Ottawa Treaty on the prohibition
of landmines as a governance benchmark, offering a historically grounded
standard against which current responses can be assessed.

The paper proceeds as follows. Section 2 reviews the relevant literature
on AI governance, Historical Institutionalism and the Ottawa Treaty
analogy. Section 3 outlines the methodology. Section 4 presents the
empirical findings, supported by charts and figures. Section 5 provides
a discussion and analysis of those findings.

\section{Literature Review}\label{literature-review}

{[}Content continues with all sections\ldots{]}

\section*{References}\label{references}
\addcontentsline{toc}{section}{References}

\phantomsection\label{refs}
\begin{CSLReferences}{1}{0}
\bibitem[\citeproctext]{ref-kostenko2022}
Kostenko, O. V. 2022. {``The Probability of Military Aggression of
Autonomous {AI}: Assumptions or Imminent Reality (Analyzing the Facts of
{Russian} War Against {Ukraine}).''} \emph{Analitichno-Porivnialne
Pravoznavstvo} 1: 179--83.
\url{https://doi.org/10.24144/2788-6018.2022.01.33}.

\bibitem[\citeproctext]{ref-RakhmetovMurzagulova2025}
Rakhmetov, Baurzhan, and Kamila Murzagulova. 2025. {``Artificial
Intelligence in Warfare: The Case of the Russia-Ukraine War.''}
\emph{Journal of Strategic Security} 18 (4): 64--77.
\url{https://doi.org/10.5038/1944-0472.18.4.2470}.

\bibitem[\citeproctext]{ref-sweijs2024}
Sweijs, Tim, and Stefanie Romansky. 2024. {``International Norms
Development and {AI} in the Military Domain.''} Centre for International
Governance Innovation \& HCSS.
\url{https://www.cigionline.org/publications/international-norms-development-and-ai-in-the-military-domain/}.

\bibitem[\citeproctext]{ref-tallberg2023}
Tallberg, Jonas, Eva Erman, Markus Furendal, Jakob Geith, Mark Klamberg,
and Mattias Lundgren. 2023. {``The Global Governance of Artificial
Intelligence: Next Steps for Empirical and Normative Research.''}
\emph{International Studies Review} 25 (3).
\url{https://doi.org/10.1093/isr/viad040}.

\end{CSLReferences}



\end{document}
